\documentclass[10pt]{article}
\usepackage[margin=0.5in]{geometry}
\usepackage{amsmath,amssymb,multicol,microtype}
\usepackage[T1]{fontenc}
\usepackage{lmodern}
\pagestyle{empty}
\setlength{\parindent}{0pt}
\setlength{\parskip}{4pt}
\setlength{\columnsep}{22pt}
\setlength{\abovedisplayskip}{5pt}
\setlength{\belowdisplayskip}{5pt}
\newcommand{\T}{\mathbb T}
\newcommand{\R}{\mathcal R}
\begin{document}
\begin{center}
{\Large\bfseries A four-phase, positive-sheet reconstruction}\\[3pt]
{\normalsize Complex Hadamard matrices of order six}
\end{center}
\textbf{Scope.} The map below covers, up to row/column permutations and unit-phase scalings, \textbf{every order-six Hadamard class except Tao and $F_6$}. Its domain is precisely
$\{(a,b,c,d)\in\T^4:\R>0\}$, where $\T=\{z:|z|=1\}$.
It is a \textbf{surjective covering, not a one-to-one parametrization}: different corners can represent the same class. Matrices are unnormalized: $|H_{ij}|=1$ and $HH^\dagger=6I_6$.

\begin{multicols}{2}
\textbf{1. Seeds and shared radical.} Put
\[
E=\begin{pmatrix}1&1&1\\1&a&b\\1&c&d\end{pmatrix},\qquad
G_h=6I_3-EE^\dagger.
\]
The horizontal and vertical data are
\[
\begin{aligned}
s_h&=-(1+a+b),&t_h&=-(1+c+d),\\
r_h&=-(1+c/a+d/b),\\
s_v&=-(1+a+c),&t_v&=-(1+b+d),\\
r_v&=-(1+b/a+d/c).
\end{aligned}
\]
Using $(s,t,r)=(s_h,t_h,r_h)$, define
\[
\begin{aligned}
X&=|s|^2,\quad Y=|t|^2,\quad Z=|r|^2,\\
p&=X+Y+Z,\quad J=2\Re(sr\bar t),\\
W&=XY+XZ+YZ-XYZ,\\
g&=27-3p+J,\quad \eta=9-p+J,\\
\chi&=W-(p-9)(p-18)+(p-15)J,\\
\R&=4g^2\eta-\chi^2>0,\\
z_h&=\frac{\chi+i\sqrt{\R}}{2g},\qquad z_v=\bar z_h.
\end{aligned}
\]
The same $g,\eta,\chi,\R$ result from the vertical data. Take $\sqrt{\R}>0$.

\textbf{2. Side recipe.} Evaluate the following twice: at $(s_h,t_h,r_h,z_h)$ and at $(s_v,t_v,r_v,z_v)$.
\[
\begin{aligned}
P_0&=rs^2+\bar r t^2+(3-|r|^2)st,\\
P_1&=t\bar s^2+r^2\bar t+(3-|t|^2)r\bar s,\\
P_2&=s\bar t^2+\bar r^2\bar s+(3-|s|^2)\bar r\bar t,\\
(c_0,c_1,c_2)&=(9-|r|^2,9-|t|^2,9-|s|^2),\\
\lambda_j&=\frac{P_j(g+c_j\bar z)}{|P_j|^2}\quad(j=0,1,2).
\end{aligned}
\]
\columnbreak
Continue the same recipe with
\[
\begin{aligned}
u&=\lambda_0\bar\lambda_1,\quad \mu=s/3,\quad\nu=t/3,\\
n&=1-|\mu|^2,\qquad k=r/3-\nu\bar\mu,\\
\alpha&=\mu^2-u\bar\mu,\quad\gamma=2\mu\nu-z\lambda_0/3,\\
\delta&=\sqrt{\frac{g}{27}\lambda_0\bar\lambda_1\bar\lambda_2},\\
d_F&=\frac{27}{g}(n\gamma-2\alpha k)\bar\delta.
\end{aligned}
\]
The number $d_F$ is real. Choose the sign of $\delta$ so that $d_F\geq0$; either sign is allowed at equality. Set
\[
\begin{aligned}
\rho&=(n+d_F)/2,\qquad q=\alpha/\rho,\\
c_+&=(k-\delta\bar q)/n,\quad c_-=(kq+\delta)/n,\\
\sigma&=u+\mu q,\qquad \xi=\sigma^2/\rho,\\
M&=\begin{pmatrix}
1&0&0\\\mu&\sigma&q\bar\sigma\\\nu&c_+\sigma&c_-\bar\sigma
\end{pmatrix},\qquad w=\sqrt[3]{\xi}.
\end{aligned}
\]
This gives $(M_h,w_h)$ and $(M_v,w_v)$. Here $|\xi|=1$; any cube root may be used. Its three choices only cycle the corresponding last three rows or columns. The square-root tie is also a side-permutation choice.

\textbf{3. Assemble $H_+(a,b,c,d)$.} Define
\[
\omega=e^{2\pi i/3},\quad
F_3=\begin{pmatrix}1&1&1\\1&\omega&\omega^2\\1&\omega^2&\omega\end{pmatrix},
\]
\[
\begin{aligned}
P(w)&=\operatorname{diag}(1,\bar w,w),\\
B&=M_hP(w_h)F_3,\\
C&=F_3P(w_v)M_v^T,\\
D&=-\frac1g\,CE^\dagger\operatorname{adj}(G_h)B.
\end{aligned}
\]
\[
\boxed{\displaystyle H_+(a,b,c,d)=\begin{pmatrix}E&B\\C&D\end{pmatrix}.}
\]
The adjugate is the transpose of the cofactor matrix: its entries are signed $2\times2$ determinants. Thus this is an explicit radical formula entirely in the four seeds. The symbol $\dagger$ denotes conjugate transpose; $T$ denotes transpose.
\end{multicols}
\vspace{-3pt}
\textbf{Validity and coverage.} On $\R>0$, $g>0$, all $P_j\ne0$ and $n,\rho>0$; no additional guards or flatness test are needed for this matched pairing. A negative-sheet representative can be moved to its opposite corner and dephased to reverse the sheet while preserving $\R>0$. Hence one positive pairing suffices globally, although it does not enumerate every completion of a prescribed corner. There are three square roots and two phase cube roots per reconstruction, with no further continuous parameters.

{\footnotesize\textbf{External proof inputs.} Strict-corner coverage, matched-completion exhaustiveness (including the exceptional Tao pairing), and opposite-corner sheet reversal are taken as proved elsewhere by the collaborator; they are not proved or independently certified in this note. Sources: \textit{Comprehensive Handoff}, \S9 (21 September 2026); \textit{Single-Sheet Conjugate Fourier Form}, \S6; and \textit{Unitary Core Phase Gate Form}.}
\end{document}
